% This is the main file for the document. It contains the preamble and the document body.，
% The preamble contains the document class, the title, the author, and any packages that are needed.
% 如果需要中文支持,请将 \documentclass{MathNote} 替换为 \documentclass{MathNoteCN}

\documentclass{MathNoteCN}

\title{数学竞赛研讨课第一节}
\author{黄宇辰}

\begin{document}
	\maketitle
	\section{一些常用不等式}
	\begin{theorem}{均值不等式}{}
		对于任意的正数 $ a_1, a_2, \ldots, a_n $，都有
		\[\frac{a_1 + a_2 + \cdots + a_n}{n} \geq \sqrt[n]{a_1 a_2 \cdots a_n}.\]
	\end{theorem}
	\begin{proof}
		(1)向前数学归纳法: \\
		归纳假设: \\
		设 $ P(n) $ 为命题“对于任意的正数 $ a_1, a_2, \ldots, a_n $，都有
		\[\frac{a_1 + a_2 + \cdots + a_n}{n} \geq \sqrt[n]{a_1 a_2 \cdots a_n}.\]
		\textbf{基例}: \\
		$ n=1 $ 时，显然成立。\\
		\textbf{假设}: \\
		假设 $ P(2^{k}) $ 成立，即对于任意的正数 $ a_1, a_2, \ldots, a_{2^k} $，都有
		\[\frac{a_1 + a_2 + \cdots + a_{2^k}}{2^k} \geq \sqrt[2^k]{a_1 a_2 \cdots a_{2^k}}.\]

		那么对于 $ n=2^{k+1} $，我们有
		\[\frac{a_1 + a_2 + \cdots + a_{2^k} + a_{2^k+1}}{2^{k+1}} =\frac{a_1+a_2+\cdots + a_{2^k}+a_{2^k+1}+}{2}\frac{a_1 + a_2 + \cdots + a_{2^k} + a_{2^k+1}}{2^k} \geq \sqrt[2^{k+1}]{a_1 a_2 \cdots a_{2^k} a_{2^k+1}}.\]

		因此，命题 $ P(n) $ 对任意正整数 $ n $ 成立。
	\end{proof}
	But this immediately implies the following corollary:
	\begin{corollary}{Riemann's}{}
		Every non-trivial zero of the Riemann $ \zeta $ function has real part one-half.
	\end{corollary}
	Which we can demonstrate with an example:
	\begin{example}{Poincare}{}
		Consider a simply connected, closed 3-manifold. Notice that it is homeomorphic to the 3-sphere!
	\end{example}
	\begin{lemma}
		This is a lemma.
	\end{lemma}
	\begin{proposition}
		This is a proposition.
	\end{proposition}
	\begin{note}
		This is a note.
	\end{note}
		\begin{definition}{Limit of a function}{}
	If, for every $ \epsilon >0 $ there exists some $ \delta >0 $ such $ 0<x-a<\delta $ implies $ |f(x)-L|<\epsilon $ then we say that the function $ f $ has a limit of $ L $ at $ a $ and we write
	\[\lim_{x\to a}f(x)=L.\]
	\end{definition}
\end{document}
